\subsubsection*{命题 3.2.}
\underline{由满足 $X\cap H$ 恰有一个奇点且该点为普通二次奇点的超平面 $H$ 构成的子集 $B\subset\check X$ 是开的.}

这是基于 R. Elkik 的“Hensel 模空间”理论(XV, 1.1.4)的 XV, 1.3.2 的特例. 若只考虑\underline{非退化}二次点, 即 $\operatorname{char}(k)\ne2$ \underline{或} $n=\dim X$ 为偶数, 则也可立即由下面的 3.3 得出; 后者的证明不使用 3.2.

\subsubsection*{推论 3.2.1.}
\underline{在 3.2 的假设下, 设 $F=\check X-B$, 它是 $\check X$ 的闭子集. 则 Lefschetz 线束构成 $\operatorname{Gr}(1,\check{\mathbb P}^{r})$ 的开子集, 且以下条件等价:}
\begin{enumerate}
\item[3.2.2.] \underline{给定嵌入 $i:X\hookrightarrow\mathbb P^r$ 为 Lefschetz 嵌入;}
\item[3.2.3.] \underline{对给定嵌入存在一个超平面 Lefschetz 线束;}
\item[3.2.4.] \underline{$F$ 在 $\check{\mathbb P}^{r}$ 中的余维至少为二.}
\end{enumerate}

\noindent\underline{证明.}
直线 $D\subset\check{\mathbb P}^{r}$ 为 Lefschetz 线束, 当且仅当:
\begin{enumerate}
\item[(3.2.5)] $D$ 的轴与 $X$ 横截相交;
\item[(3.2.6)] $D$ 不包含于 $\check X$;
\item[(3.2.7)] $D$ 不与 $F$ 相交.
\end{enumerate}
条件 (3.2.5) 及 (3.2.6) 定义 $\operatorname{Gr}(1,\check{\mathbb P}^{r})$ 的开子集; 由于 $F$ 闭, (3.2.7) 也定义开子集. 因而 Lefschetz 线束构成开子集.

由于 $\dim\check X\leq r-1$, (3.2.6) 中的开子集非空, 从而稠密. (3.2.5) 中的开子集同样非空且稠密. 事实上, 选择一个与 $X$ 横截相交的超平面 $H$; 由于 $\dim\check X\leq r-1$, 这样的 $H$ 存在. 因为 $\dim\check{(X\cap H)}\leq r-1$, 存在一个与 $X\cap H$ 横截相交的超平面 $G$, 于是经过 $H$ 和 $G$ 的直线 $D$ 的轴与 $X$ 横截相交.\EditorialFootnote{维数界核对: 此处三次使用 3.1.4 的对偶簇维数界, 所需对象分别为 $\check X$、$\check X$ 和 $\check{(X\cap H)}$, 而非原簇本身. 法文原书第 220 页第二、第三处可见对偶记号, 第三处括号内又印为 $X\cap G$, 与紧后的 $X\cap H$ 不一致; 英文底本漏去对偶记号, 并已将第三处改作 $X\cap H$. 本译文按所引用的维数界和论证顺序明示这些对象.} 因此此嵌入为 Lefschetz 嵌入, 当且仅当 (3.2.7) 的开子集非空, 从而稠密. 由于 $F$ 闭, 此条件成立当且仅当 $F$ 在 $\check{\mathbb P}^{r}$ 中的余维至少为二(参见 EGA V, 或 Mumford, \emph{Introduction to Algebraic Geometry}, 第~88~页, 推论~3).

\subsubsection*{推论 3.2.8.}
\underline{在 3.2 的假设下, 假设 $i:X\hookrightarrow\mathbb P^r$ 为 Lefschetz 嵌入. 设 $H_0$ 为一个与 $X$ 横截相交的超平面, $\operatorname{Gr}(1,\check{\mathbb P}^{r})_{H_0}$ 为 $\operatorname{Gr}(1,\check{\mathbb P}^{r})$ 中由过 $H_0$ 的直线组成的子簇. 则过 $H_0$ 的 Lefschetz 线束构成 $\operatorname{Gr}(1,\check{\mathbb P}^{r})_{H_0}$ 的稠密开子集.}

\noindent\underline{证明.}
由 3.2.1, 这些束构成 $\operatorname{Gr}(1,\check{\mathbb P}^{r})_{H_0}$ 的开子集. 刚才用于证明开子集 (3.2.5) 和 (3.2.6) 非空的论证也表明, 它们与此子簇的交非空. 尚需证明存在某些过 $H_0$ 的直线不与 $F$ 相交.

选取不含 $H_0$ 的超平面 $T\subset\check{\mathbb P}^{r}$. 有同构
\begin{center}
\begin{tikzcd}[column sep=huge]
\operatorname{Gr}(1,\check{\mathbb P}^{r})_{H_0}
  & T \arrow[l,"\sim"']\\
\substack{\text{经过下列两点的直线}\\ H_0\text{ 和 }\tau}
  & \tau \arrow[l,mapsto]
\end{tikzcd}
\end{center}
因此 $\operatorname{Gr}(1,\check{\mathbb P}^{r})_{H_0}$ “就是”一个 $r-1$ 维射影空间. 其中经过 $F$ 中一点的直线构成一个维数至多 $r-2$ 的闭子集, 因为它是态射
\begin{center}
\begin{tikzcd}[column sep=huge]
F \arrow[r] & \operatorname{Gr}(1,\check{\mathbb P}^{r})_{H_0}\\
f \arrow[r,mapsto] &
  \substack{\text{经过 }H_0\\\text{及经过 }f}.
\end{tikzcd}
\end{center}
的像.

\subsubsection*{命题 3.3.}
\underline{态射 $\varphi:\mathbb P(N)\to\check{\mathbb P}^{r}$ 在闭点 $(x_0,H)$ 处非分歧, 当且仅当 $x_0$ 为 $X\cap H$ 的非退化二次奇点(1.2). 特别地, $\mathbb P(N)$ 中由满足 $x$ 为 $X\cap H$ 的非退化二次奇点的点 $(x,H)$ 组成的子集为开集.}